\section*{Latihan dan Refleksi}
\addcontentsline{toc}{section}{Latihan dan Refleksi}

\subsection*{Latihan}
\begin{enumerate}[label=\arabic*.,leftmargin=*]
  \item Jelaskan perbedaan token, lexeme, dan pattern dengan contoh.
  \item Tuliskan pattern untuk identifier yang boleh mengandung garis
        bawah dan angka, tetapi tidak boleh diawali angka.
  \item Tuliskan pattern untuk bilangan pecahan tanpa bagian
        eksponen.
  \item Mengapa operator dua karakter harus ditempatkan sebelum
        operator satu karakter?
  \item Berikan contoh di mana maximal munch mencegah kesalahan
        parsing.
  \item Rancang aturan untuk komentar blok yang tidak bersarang.
  \item Jelaskan perilaku lexer ketika menemukan string yang tidak
        ditutup.
  \item Bandingkan strategi kata kunci sebagai identifier yang
        diperiksa tabel dengan strategi kata kunci sebagai pattern
        tersendiri.
\end{enumerate}

\subsection*{Refleksi}
\begin{reflectionbox}
\begin{enumerate}[label=\arabic*.,leftmargin=*]
  \item Bagian mana dari spesifikasi token yang paling sulit Anda
        tulis? Mengapa?
  \item Apakah Anda dapat menjelaskan mengapa lexer tidak perlu
        mengetahui tipe variabel?
  \item Kasus batas apa yang sering Anda lupakan saat menulis
        spesifikasi?
\end{enumerate}
\end{reflectionbox}
